\documentclass[FM,DP]{tulthesis}

\usepackage{polyglossia}
\setdefaultlanguage{czech}
\usepackage{fontspec}
\usepackage{xunicode}
\usepackage{xltxtra}
\setsansfont[Mapping=tex-text,BoldFont={* Bold},Numbers=OldStyle]{Myriad Pro}
\usepackage{hyperref}
\hypersetup{colorlinks=true, linkcolor=tul, urlcolor=tul, citecolor=tul}
\usepackage{graphicx}
\usepackage{listings}
\usepackage[toc,page]{appendix}
\usepackage{amsmath}
\usepackage{amssymb}
\newcommand{\argument}[1]{{\ttfamily\color{\tulcolor}#1}}
\newcommand{\prikaz}[1]{\argument{\textbackslash #1}}
\newenvironment{myquote}{\begin{list}{}{\setlength\leftmargin\parindent}\item[]}{\end{list}}
\newenvironment{listing}{\begin{myquote}\color{\tulcolor}}{\end{myquote}}
\sloppy


%%%%%%%%%%%%%%%%%%%%%%%%%%%%
\TULtitle{Vyhledávání jako služba}{Search as a Service}
\TULprogramme{Aplikovaná informatika}{Applied Informatics}
\TULbranch{Informační systémy a technologie}{Information Technologies}
\TULauthor{Bc. Luděk Veselý}
\TULsupervisor{Prof. Ing. Zdeněk Molnár, CSc.}
\TULyear{2017}


%%%%%%%%%%%%%%%%%%%%%%%%%%%%
\begin{document}
\ThesisStart{male}


%%%%%%%%%%%%%%%%%%%%%%%%%%%%
\begin{abstractCZ}
Tato diplomová práce popisuje návrh a tvorbu fulltextového vyhledávání poskytovaného jako služba.
\end{abstractCZ}

\begin{klicovaslovaCZ}
Fulltext, Elasticsearch
\end{klicovaslovaCZ}

\vspace{2cm}

\begin{abstractEN}
This diploma thesis describes creation of fulltext search service.
\end{abstractEN}

\begin{klicovaslovaEN}
Fulltext, Elasticsearch
\end{klicovaslovaEN}


%%%%%%%%%%%%%%%%%%%%%%%%%%%%
\begin{acknowledgement}
Rád bych poděkoval Ing. Ivanovi Jelínkovi za rady a pomoc při řešení.
\end{acknowledgement}


%%%%%%%%%%%%%%%%%%%%%%%%%%%%
\tableofcontents
\clearpage


%%%%%%%%%%%%%%%%%%%%%%%%%%%%
%\begin{abbrList}
%\textbf{ES} & Elasticsearch\\
%\textbf{DBS} & Databázový systém\\
%\end{abbrList}


%%%%%%%%%%%%%%%%%%%%%%%%%%%%
\chapter{Úvod}

V~této diplomové práci práci...

\section{Cíle práce}
...

\section{Cílová skupina}
...

\section{Použité metody}
...

\section{Struktura práce}
...


%%%%%%%%%%%%%%%%%%%%%%%%%%%%
\chapter{Komentovaná rešerše informačních zdrojů}
...


%%%%%%%%%%%%%%%%%%%%%%%%%%%%
\chapter{Analýza byznys požadavků}
- Vysvetleni problemu, ktery je resen
- Specifika prostredi e-commerce, definice kontextu, ve kterem bude nastroj pouzivan
- Definice pozadavku na samotne vyhledavani, akceptacni kriteria
- Technicke pozadavky (datove formaty, rychlost odezvy...)
- Popis oboru, kterych se prace dotyka (ziskavani informaci, NLP, big data...)

\chapter{Teorie fulltextového vyhledávání}
- Popis problemu indexovani dokumentu
- Analyza dat pouzitych k vyhledavani
- Zpracovani prirozeneho jazyka
  - Sklonovani, casovani, tvaroslovi obecne
  - Vaha slov, stop slova
  - Preklepy
  - Synonyma, zargon
  - Lemmatizace
  - Tezaurus

\chapter{Návrh řešení}

\section{Nastroje pouzitelne pro fulltextove vyhledavani}
- Popis a porovnani vhodnych nastroju
  - MySQL
  - Elasticsearch
  - Sphinx
  - Algolia
  - Navrh architektury aplikace

\section{Návrh aplikace}
- Definice pripadu uziti
- Navrh architektury aplikace
- Vyber technologii a nastroju (jazyk, framework...)

\chapter{Implementace}

\section{Architektura aplikace}
- Frontend vs backend (API)
- Nastaveni a nasazeni Elasticsearch
  - Implementace indexace
  - Implementace vyhledavani
- Backend: popis trid, rozhrani, implementacni detaily
- Frontend: UI, komunikace s backendem
- Deployment, CI

\section{Testovani, overeni}
- Overeni funkcnosti a prinosu
- Testovani, zda jsou splnena akceptacni kriteria
- Porovnani s stavajicim resenim vyhledavani na konkretnim priklade


%%%%%%%%%%%%%%%%%%%%%%%%%%%%
\chapter{Závěr}

V této diplomové práci jsem...
Diskuse mozneho budouciho rozsirovani

\section{Dosažení vytčených dílů}
...


%%%%%%%%%%%%%%%%%%%%%%%%%%%%
\begin{thebibliography}{Mm99}

\bibitem{strossa} STROSSA, Petr. \emph{Počítačové zpracování přirozeného jazyka}. \\
Vyd.~1. Praha: Oeconomica, 2011 316~s. ISBN 978-80-245-1777-3.

\bibitem{searching} AYSE Göker a DAVIES John. \emph{Information Retrieval: Searching in the 21st Century}. \\
Vyd.~1. Library of Congress Cataloging-in-Publication Data, 2009 295~s. ISBN: 978-0-470-02762-2.

\bibitem{mining} AGGARWAL Charu C., ZHAI ChengXiang. \emph{Mining Text Data}. \\
Springer New York Dordrecht Heidelberg London, 2000 522~s. ISBN 978-1-4614-3222-7.

\bibitem{es-guide} CLINTON Gormley, ZACHARY Tong. \emph{Elasticsearch: The Definitive Guide}. \\
O'Reilly Media, 2015 724~s. ISBN 978-1-4493-5854-9.

\bibitem{go-in-action} KENNEDY William. \emph{Go in Action}. \\
Manning Publications Co, 2016 264~s. ISBN 978-1-6172-9178-4.

\bibitem{elastic-reference} Elasticsearch. \emph{Elasticsearch Reference} [online]. \\
2017 [cit. 2017-03-03]. Dostupné z:\\
\url{https://www.elastic.co/guide/en/elasticsearch/reference/current/index.html}

\end{thebibliography}


%%%%%%%%%%%%%%%%%%%%%%%%%%%%
\chapter*{Rejstřík}
...


%%%%%%%%%%%%%%%%%%%%%%%%%%%%
\appendix

\chapter{Obsah přiloženého DVD}

\begin{itemize}
\item Soubor Diplomova\_prace\_2017\_Ludek\_Vesely.pdf
\begin{itemize}
	\item Text diplomové práce
\end{itemize}

\end{itemize}


%\chapter{Graf jako příloha}
%\begin{figure}[h]
%\center
%\includegraphics[width=0.86\textwidth]{todo.pdf}
%\end{figure}


%%%%%%%%%%%%%%%%%%%%%%%%%%%%
\end{document}
